\documentclass[12pt,a4paper]{article}

% ── Packages ─────────────────────────────────────────────────────────────────
\usepackage[utf8]{inputenc}
\usepackage[T1]{fontenc}
\usepackage[turkish]{babel}
\usepackage{geometry}
\usepackage{graphicx}
\usepackage{xcolor}
\usepackage{titlesec}
\usepackage{booktabs}
\usepackage{array}
\usepackage{longtable}
\usepackage{enumitem}
\usepackage{hyperref}
\usepackage{fancyhdr}
\usepackage{mdframed}
\usepackage{tikz}
\usepackage{pgfplots}
\usepackage{multicol}
\usepackage{amsmath}
\usepackage{microtype}
\usepackage{setspace}
\usepackage{fontawesome5}

\pgfplotsset{compat=1.18}

% ── Sayfa Düzeni ─────────────────────────────────────────────────────────────
\geometry{
  top=2.5cm, bottom=2.5cm,
  left=2.8cm, right=2.8cm
}

% ── Renkler ──────────────────────────────────────────────────────────────────
\definecolor{primary}{HTML}{1E40AF}      % indigo-700
\definecolor{accent}{HTML}{3B82F6}       % blue-500
\definecolor{dark}{HTML}{111827}         % gray-900
\definecolor{muted}{HTML}{6B7280}        % gray-500
\definecolor{light}{HTML}{EFF6FF}        % blue-50
\definecolor{success}{HTML}{16A34A}      % green-600
\definecolor{danger}{HTML}{DC2626}       % red-600
\definecolor{warning}{HTML}{D97706}      % amber-600
\definecolor{sectionbg}{HTML}{F0F4FF}

% ── Hyperlink Ayarları ────────────────────────────────────────────────────────
\hypersetup{
  colorlinks=true,
  linkcolor=primary,
  urlcolor=accent,
  citecolor=primary,
  pdfauthor={Lagent Ekibi},
  pdftitle={Lagent — Yatırımcı Sunumu},
  pdfsubject={Hukuki Sözleşme Analiz Platformu},
}

% ── Başlık Stilleri ───────────────────────────────────────────────────────────
\titleformat{\section}
  {\color{primary}\Large\bfseries}
  {}
  {0em}
  {}
  [\color{accent}\titlerule]

\titleformat{\subsection}
  {\color{dark}\large\bfseries}
  {}
  {0em}
  {}

\titleformat{\subsubsection}
  {\color{primary}\normalsize\bfseries}
  {}
  {0em}
  {}

\titlespacing*{\section}{0pt}{1.8em}{0.8em}
\titlespacing*{\subsection}{0pt}{1.2em}{0.5em}

% ── Üst/Alt Bilgi ─────────────────────────────────────────────────────────────
\pagestyle{fancy}
\fancyhf{}
\fancyhead[L]{\small\color{muted}\textbf{Lagent} — Gizli Yatırımcı Sunumu}
\fancyhead[R]{\small\color{muted}Mayıs 2026}
\fancyfoot[C]{\small\color{muted}\thepage}
\renewcommand{\headrulewidth}{0.3pt}

% ── Kutucuk Stilleri ──────────────────────────────────────────────────────────
\newmdenv[
  backgroundcolor=light,
  linecolor=accent,
  linewidth=1.5pt,
  leftline=true, topline=false, rightline=false, bottomline=false,
  innerleftmargin=12pt,
  innerrightmargin=12pt,
  innertopmargin=10pt,
  innerbottommargin=10pt,
  skipabove=10pt,
  skipbelow=10pt,
]{highlightbox}

\newmdenv[
  backgroundcolor=sectionbg,
  linecolor=primary,
  linewidth=0pt,
  roundcorner=6pt,
  innerleftmargin=14pt,
  innerrightmargin=14pt,
  innertopmargin=12pt,
  innerbottommargin=12pt,
  skipabove=8pt,
  skipbelow=8pt,
]{statbox}

% ═════════════════════════════════════════════════════════════════════════════
\begin{document}
% ═════════════════════════════════════════════════════════════════════════════

% ── Kapak Sayfası ─────────────────────────────────────────────────────────────
\begin{titlepage}
  \thispagestyle{empty}
  \begin{tikzpicture}[remember picture, overlay]
    \fill[primary] (current page.north west) rectangle
      ([yshift=-5cm]current page.north east);
    \fill[accent!30!primary] (current page.south west) rectangle
      ([yshift=1.5cm]current page.south east);
  \end{tikzpicture}

  \vspace*{1cm}
  \begin{center}
    {\color{white}\fontsize{48}{56}\selectfont\textbf{Lagent}}\\[0.3em]
    {\color{white!80!primary}\Large Yapay Zeka Destekli Hukuki Sözleşme Analiz Platformu}\\[3.5cm]

    {\color{dark}\fontsize{22}{28}\selectfont\textbf{Yatırımcı Sunumu}}\\[0.5em]
    {\color{muted}\large Seri A Öncesi Tur}\\[0.5em]
    {\color{muted}\large Mayıs 2026}

    \vfill
    \begin{mdframed}[backgroundcolor=light,linecolor=accent,linewidth=1pt,roundcorner=6pt]
      \centering
      {\small\color{muted}\textbf{GİZLİLİK NOTU}}\\[3pt]
      {\small\color{dark}Bu belge yalnızca yetkili alıcılara yönelik olup gizlidir.
      İzinsiz çoğaltılması, dağıtılması veya üçüncü şahıslarla paylaşılması yasaktır.}
    \end{mdframed}
    \vspace{1cm}
    {\color{muted}\normalsize Mustafa Emir Ceran, Mert Ayrancı, Osman Gazi Atalay}
  \end{center}
\end{titlepage}

\newpage
\setcounter{page}{1}
\onehalfspacing

% ═════════════════════════════════════════════════════════════════════════════
\section*{İçindekiler}
% ═════════════════════════════════════════════════════════════════════════════
\begin{multicols}{2}
\begin{enumerate}[leftmargin=*, label=\textbf{\arabic*.}]
  \item Yönetici Özeti
  \item Problem: Sözleşme Bunalımı
  \item Çözüm: Lagent Platformu
  \item Ürün Özellikleri
  \item Teknoloji Mimarisi
  \item Pazar Fırsatı
  \item İş Modeli
  \item Rekabet Analizi
  \item Yol Haritası
  \item Ekip
  \item Finansal Projeksiyonlar
  \item Yatırım Talebi
\end{enumerate}
\end{multicols}

\newpage

% ═════════════════════════════════════════════════════════════════════════════
\section{Yönetici Özeti}
% ═════════════════════════════════════════════════════════════════════════════

\begin{highlightbox}
\textbf{\color{primary}Tek Cümle:} Lagent, bireyler ve KOBİ'lerin hukuki sözleşmelerini dakikalar
içinde yapay zeka ile analiz etmelerini, risklerini anlamalarını ve güvenle müzakere etmelerini
sağlayan çok-ajanlı bir SaaS platformudur.
\end{highlightbox}

\vspace{0.5em}

Türkiye'de her yıl \textbf{500 milyonun üzerinde sözleşme} imzalanmaktadır. Bu sözleşmelerin büyük
çoğunluğu kira, tedarik, istihdam ve hizmet anlaşmalarından oluşmakta; ancak imzalayanların
\textbf{\%78'i} sözleşmeyi tam olarak anlayamadan imzaladığını kabul etmektedir. Hukuki danışmanlık
masrafları ise ortalama \textbf{sözleşme başına 1.500–5.000 TL} düzeyindedir.

Lagent bu boşluğu kapatır. Kullanıcılar sözleşmelerini platforma yükler; yapay zeka destekli
çok-ajanlı sistem her maddeyi ayrıştırır, risklerini puanlar, politika uyumunu kontrol eder ve
revizyon önerileri sunar. \textbf{Tüm bu süreç 2 dakikadan kısadır.}

\vspace{0.8em}

\begin{statbox}
\begin{center}
\begin{tabular}{p{3.5cm}p{3.5cm}p{3.5cm}p{3.5cm}}
\centering\textbf{\color{primary}\Large 500M+} &
\centering\textbf{\color{primary}\Large \%78} &
\centering\textbf{\color{primary}\Large 2 dk} &
\centering\textbf{\color{primary}\Large \%80} \\[2pt]
\centering{\small Yıllık sözleşme\newline (Türkiye)} &
\centering{\small Anlamadan\newline imzalayanlar} &
\centering{\small Analiz süresi\newline (10 sayfa)} &
\centering{\small Maliyet azalma\newline potansiyeli}
\end{tabular}
\end{center}
\end{statbox}

% ═════════════════════════════════════════════════════════════════════════════
\section{Problem: Sözleşme Bunalımı}
% ═════════════════════════════════════════════════════════════════════════════

\subsection{Kimin Sorunu?}

\textbf{Bireyler} kira sözleşmelerini, iş akitlerini ve hizmet anlaşmalarını avukata başvurmadan
imzalar. Karmaşık hukuki dil, gizlenmiş cezai şartlar ve eksik hükümler fark edilmez. Sonuç:
tahmin edilemeyen mali ve hukuki riskler.

\textbf{KOBİ'ler} tedarikçi, ortak ve müşteri sözleşmelerini sınırlı hukuki bütçeyle yönetmek
zorundadır. İnsan kaynağı yetersiz, profesyonel hukuki inceleme ise pahalı ve yavaştır.

\textbf{Hukuk Profesyonelleri} standart madde kontrolüne zaman kaybeder; karmaşık müzakere ve
stratejik analiz için yeterli alan bulamazlar.

\subsection{Sorunun Büyüklüğü}

\begin{itemize}[leftmargin=*, itemsep=4pt]
  \item \textbf{Zaman maliyeti:} Bir avukat ortalama 10 sayfalık sözleşmeyi 2–4 saatte inceler.
  \item \textbf{Parasal maliyet:} Türkiye'de sözleşme başına hukuki danışmanlık 1.500–5.000 TL;
  yıllık hukuk danışmanlık harcamaları KOBİ'lerde ortalama 50.000–200.000 TL.
  \item \textbf{Erişim eşitsizliği:} Hukuki hizmetler büyük şirket ve varlıklı bireyler için
  erişilebilir; KOBİ ve bireyler için değil.
  \item \textbf{İnsan hatası:} Manuel incelemede kritik maddelerin gözden kaçma oranı \%23
  (Harvard Law Review, 2023).
  \item \textbf{Dil engeli:} Hukuki Türkçe, ortalama kullanıcı tarafından anlaşılması güç teknik
  terminoloji içerir.
\end{itemize}

\begin{highlightbox}
\textbf{\color{danger}Sonuç:} Her yıl Türkiye'de sözleşme anlaşmazlıklarından kaynaklanan
ekonomik kayıp tahminen \textbf{15–25 milyar TL} civarındadır. Ticari mahkemelerdeki davaların
\textbf{\%40'ı} sözleşme belirsizliklerinden kaynaklanmaktadır.
\end{highlightbox}

% ═════════════════════════════════════════════════════════════════════════════
\section{Çözüm: Lagent Platformu}
% ═════════════════════════════════════════════════════════════════════════════

\subsection{Ürün Vizyonu}

Lagent, hukuki sözleşme incelemesini \textbf{saniyeler içinde, demokratik biçimde} herkes için
erişilebilir kılan, yapay zeka destekli çok-ajanlı bir platformdur. Sistem avukatın yerini almaz;
\textbf{en güvenilir hukuki asistanı} gibi davranır: analiz eder, uyarır, önerir — nihai kararı
her zaman kullanıcı verir.

\subsection{Nasıl Çalışır?}

\begin{enumerate}[leftmargin=*, itemsep=6pt]
  \item \textbf{Yükle:} Kullanıcı PDF veya DOCX formatındaki sözleşmeyi yükler ($\leq$10 MB).
  Sistem metni otomatik çıkarır; taranan belgeler için OCR devreye girer.

  \item \textbf{Analiz Et (Madde Ajanı):} Çok-ajanlı orkestratör sözleşmeyi bağımsız maddelere
  böler, her birini hukuki kategorisine (gizlilik, fesih, cezai şart, ödeme koşulları vb.) göre
  sınıflandırır.

  \item \textbf{Risk Puanla (Risk Ajanı):} Her madde kullanıcı tanımlı Playbook kurallarıyla
  anlam tabanlı eşleştirme (RAG) ve LLM değerlendirmesiyle karşılaştırılır; ticari ve hukuki
  risk puanı (0.0–1.0) üretilir. Eksik hükümler tespit edilir.

  \item \textbf{Revizyon Öner (Müzakere Ajanı):} Orta ve yüksek riskli maddelere alternatif
  metin önerilir. Kullanıcı orijinal ve önerilen metni yan yana (redline) karşılaştırır.

  \item \textbf{Onayla / Reddet:} Madde bazlı karar akışı; her karar denetim iziyle kayıt altına
  alınır.

  \item \textbf{Raporla:} Özet veya detaylı rapor PDF/DOCX olarak oluşturulur; revize edilmiş
  sözleşme taslağı indirilebilir.
\end{enumerate}

% ═════════════════════════════════════════════════════════════════════════════
\section{Ürün Özellikleri}
% ═════════════════════════════════════════════════════════════════════════════

\subsection{Temel Özellikler}

\begin{longtable}{>{\bfseries\color{primary}}p{4.5cm} p{9.5cm}}
\toprule
\textbf{\color{dark}Özellik} & \textbf{\color{dark}Açıklama} \\
\midrule
\endhead
Belge İşleme & PDF ve DOCX desteği; otomatik metin çıkarma; taranmış belgeler için OCR (Tesseract); 10 sayfalık sözleşme $<$30 saniye \\[4pt]
Madde Ayrıştırma & GPT-4o ile bağımsız maddelere otomatik bölme; 8+ hukuki kategori sınıflandırması; güven skoru ($\geq$0.7 eşiği) \\[4pt]
Playbook Politikaları & Sektöre özel kullanıcı tanımlı politikalar; kabul edilebilir / reddedilen / zorunlu / eşik tipi kurallar; çoklu playbook desteği \\[4pt]
Risk Değerlendirmesi & Ticari + hukuki çift boyutlu puanlama; 3 seviyeli risk (Düşük/Orta/Yüksek); gerekçe açıklaması; eksik hüküm tespiti \\[4pt]
Revizyon Önerileri & Orta/yüksek riskli maddelere LLM destekli alternatif metin; madde başına $<$5 saniye \\[4pt]
Redline Görünümü & Orijinal vs. önerilen metin yan yana renk kodlu karşılaştırma \\[4pt]
Onay İş Akışı & Madde bazlı onayla/reddet/revize; zorunlu insan kararı; tam denetim izi \\[4pt]
Raporlama & PDF/DOCX özet ve detaylı rapor; revize sözleşme taslağı çıktısı \\[4pt]
Gerçek Zamanlı Bildirim & WebSocket ile analiz ilerleme durumu canlı güncelleme \\[4pt]
Güvenlik & JWT kimlik doğrulama; bcrypt şifreleme ($\geq$12 karakter); RBAC; HTTPS; denetim kaydı \\
\bottomrule
\end{longtable}

\subsection{Playbook: Rekabet Gücünün Kaynağı}

\begin{highlightbox}
\textbf{\color{primary}Playbook nedir?} Kullanıcı veya kuruluşun kendi sözleşme politikalarını
tanımladığı kural kümesidir. \emph{"Depozito maksimum 3 aylık kira olsun"} ya da \emph{"Tahkim
maddesi zorunlu"} gibi kurallar sisteme eklenir; her analiz bu kurallara göre yapılır.
\end{highlightbox}

Playbook sistemi Lagent'i rakiplerinden ayıran temel farktır:
\begin{itemize}[itemsep=3pt]
  \item Sektör ve şirket özelinde politika tanımlanabilir
  \item Kurallar anlam tabanlı vektör arama (RAG) ile maddelere eşleştirilir
  \item Yeni kural eklendiğinde vektör indeksi otomatik güncellenir
  \item Eksik hükümleri gerçek zamanlı tespit eder
\end{itemize}

% ═════════════════════════════════════════════════════════════════════════════
\section{Teknoloji Mimarisi}
% ═════════════════════════════════════════════════════════════════════════════

\subsection{Çok-Ajanlı Orkestrasyon}

Lagent'in kalbi üç uzmanlaşmış yapay zeka ajanından oluşan bir orkestratördür:

\vspace{0.5em}
\begin{statbox}
\begin{tabular}{>{\bfseries\color{primary}}p{3.5cm} p{10cm}}
Madde Ajanı & Sözleşme metnini yapılandırılmış JSON formatında maddelere böler. GPT-4o (JSON mode) kullanır. Doğruluk: $\geq$\%70. \\[6pt]
Risk Ajanı & RAG (pgvector anlam araması) + LLM değerlendirmesi + deterministik kural motoru. Playbook uyumunu çapraz doğrular. \\[6pt]
Müzakere Ajanı & Yüksek/orta riskli maddeler için bağlam farkındalıklı revizyon üretir. Benzer onaylanmış örneklerden öğrenir. \\
\end{tabular}
\end{statbox}

\subsection{Teknoloji Yığını}

\begin{longtable}{>{\bfseries}p{3.5cm} p{4.5cm} p{5.5cm}}
\toprule
\textbf{Katman} & \textbf{Teknoloji} & \textbf{Açıklama} \\
\midrule
\endhead
\multicolumn{3}{l}{\color{primary}\textbf{Frontend}} \\[2pt]
Çerçeve & Next.js 14 + React 18 & SSR destekli, hızlı UI \\
Stil & Tailwind CSS + shadcn/ui & Modern, erişilebilir bileşenler \\
State & Zustand & Hafif global state yönetimi \\
Gerçek Zamanlı & Native WebSocket & Analiz ilerleme bildirimi \\
\midrule
\multicolumn{3}{l}{\color{primary}\textbf{Backend}} \\[2pt]
Framework & FastAPI 0.110+ & Async Python, yüksek performanslı API \\
ORM & SQLAlchemy 2.0 & Async veri erişimi \\
Kimlik Doğrulama & JWT + bcrypt & Stateless güvenli oturum \\
Belge İşleme & PyMuPDF + python-docx & PDF/DOCX metin çıkarma \\
OCR & Tesseract 5.x & Taranmış belge desteği \\
\midrule
\multicolumn{3}{l}{\color{primary}\textbf{Yapay Zeka}} \\[2pt]
LLM & GPT-4o / GPT-4o-mini & Düşük temp. (0.2) deterministik çıktı \\
Embedding & text-embedding-3-small & 1536 boyutlu anlam vektörleri \\
Vektör Araması & pgvector & PostgreSQL içi anlam benzerliği \\
Orkestrasyon & LangGraph & Durum makineli çok-ajan akışı \\
\midrule
\multicolumn{3}{l}{\color{primary}\textbf{Altyapı}} \\[2pt]
Veritabanı & PostgreSQL 16 + pgvector & İlişkisel + vektör depolama \\
Önbellek & Redis 7 & Oturum, durum, oran sınırlama \\
Nesne Depolama & MinIO (S3-uyumlu) & Sözleşme ve rapor dosyaları \\
Konteyner & Docker + Compose & Taşınabilir, ölçeklenebilir dağıtım \\
CI/CD & GitHub Actions & Otomatik test ve dağıtım \\
\bottomrule
\end{longtable}

\subsection{Neden Bu Mimari?}

\begin{itemize}[itemsep=4pt]
  \item \textbf{RAG + LLM kombinasyonu}: Salt LLM'den farklı olarak kullanıcı politikaları anlam
  aramasıyla maddelere özel bağlam oluşturur; halüsinasyon riski minimize edilir.
  \item \textbf{Deterministik kural motoru}: LLM çıktısı kural motoruyla çapraz doğrulanır;
  yüksek güvenlik gerektiren hukuki bağlamda kritik güvenilirlik sağlar.
  \item \textbf{pgvector}: Ayrı bir vektör veritabanı servisine gerek kalmadan PostgreSQL içinde
  anlam araması; operasyonel karmaşıklık azaltılır.
  \item \textbf{Stateless API}: Yatay ölçekleme kolaylığı; mikro-servis mimarisine hazır altyapı.
\end{itemize}

% ═════════════════════════════════════════════════════════════════════════════
\section{Pazar Fırsatı}
% ═════════════════════════════════════════════════════════════════════════════

\subsection{Hedef Pazar}

\subsubsection{Birincil: Türkiye Pazarı}

\begin{statbox}
\begin{tabular}{p{5cm} p{3.5cm} p{5cm}}
\textbf{Segment} & \textbf{Büyüklük} & \textbf{Lagent Değeri} \\
\midrule
KOBİ (Türkiye) & $\sim$8 milyon işletme & Hukuki maliyet \%60–80 azalma \\
Bireyler & $\sim$85 milyon nüfus & Kira/iş akdi analizi \\
Hukuk Büroları & $\sim$80.000 avukat & Verimlilik artışı $\times$3 \\
İK / Tedarik Ekipleri & 500.000+ yönetici & Standart sözleşme inceleme \\
\end{tabular}
\end{statbox}

\subsubsection{Pazar Boyutu Analizi (TAM / SAM / SOM)}

\begin{itemize}[itemsep=5pt]
  \item \textbf{TAM (Toplam Adreslenebilir Pazar):} Türkiye'de yıllık hukuki danışmanlık pazarı
  $\sim$\textbf{45 milyar TL} (2025). Küresel LegalTech pazarı 2030'a kadar \textbf{37 milyar \$}
  büyüklüğe ulaşacak (CAGR \%9.1).

  \item \textbf{SAM (Hizmet Verilebilir Pazar):} Sözleşme analizi ve yönetimi alt segmenti
  $\sim$\textbf{8 milyar TL} (Türkiye KOBİ + birey pazarı).

  \item \textbf{SOM (Elde Edilebilir Pazar):} 3 yıl içinde \textbf{50.000 abonelik} hedefiyle
  $\sim$\textbf{150 milyon TL} yıllık tekrar eden gelir (ARR).
\end{itemize}

\subsection{Pazar Büyümesi}

\begin{itemize}[itemsep=4pt]
  \item LegalTech benimsenmesi hızlanıyor: COVID-19 sonrası dijital iş süreçleri kalıcılaştı;
  sözleşme yönetimi dijitalleşmesi \%35 CAGR büyüyor (Gartner, 2024).
  \item Türkiye'de e-imza yaygınlaşması ve dijital sözleşme hacmi yıldan yıla \%40 artış gösteriyor.
  \item Yapay zeka ve doğal dil işleme teknolojilerinin maliyeti düşerken kalitesi hızla artıyor;
  Türkçe NLP yetenekleri 2023–2026 arasında dramatik iyileşme gösterdi.
  \item KVKK ve yeni ticaret mevzuatı uyum gereklilikleri sözleşme denetim talebini artırıyor.
\end{itemize}

% ═════════════════════════════════════════════════════════════════════════════
\section{İş Modeli}
% ═════════════════════════════════════════════════════════════════════════════

\subsection{Gelir Akışları}

\begin{longtable}{>{\bfseries\color{primary}}p{3cm} p{4cm} p{3cm} p{3.5cm}}
\toprule
\textbf{\color{dark}Plan} & \textbf{\color{dark}Hedef Kitle} & \textbf{\color{dark}Fiyat} & \textbf{\color{dark}Özellikler} \\
\midrule
\endhead
Freemium & Bireyler & Ücretsiz & 3 sözleşme/ay, temel analiz \\[4pt]
Starter & Serbest çalışanlar & 299 TL/ay & 20 sözleşme/ay, raporlama \\[4pt]
Pro & KOBİ'ler & 999 TL/ay & Sınırsız sözleşme, playbook, API \\[4pt]
Enterprise & Kurumsal / Hukuk bürosu & Özel fiyat & On-prem seçeneği, white-label, öncelikli destek \\[4pt]
Pay-per-Use & Yüksek hacim & 25 TL/sözleşme & API entegrasyonu, toplu analiz \\
\bottomrule
\end{longtable}

\subsection{Birim Ekonomisi (Pro Plan, 3. Yıl Tahmini)}

\begin{statbox}
\begin{tabular}{lrl}
\textbf{Müşteri Başına Aylık Gelir (ARPU)} & 999 TL & \\
\textbf{Müşteri Edinme Maliyeti (CAC)} & 1.200 TL & (dijital pazarlama + satış) \\
\textbf{Yaşam Boyu Değer (LTV)} & 18.000 TL & (24 ay ortalama abonelik) \\
\textbf{LTV / CAC Oranı} & 15x & \color{success}\textbf{(Hedef: $>$3x)} \\
\textbf{Brüt Marj} & \%72 & (LLM API + altyapı maliyeti çıkarıldıktan sonra) \\
\textbf{Geri Ödeme Süresi} & $\sim$1.5 ay & \\
\end{tabular}
\end{statbox}

\subsection{Büyüme Stratejisi}

\begin{enumerate}[leftmargin=*, itemsep=5pt]
  \item \textbf{Ürün Odaklı Büyüme (PLG):} Freemium katmanı ile organik kullanıcı edinimi; viral
  döngüler (sözleşme paylaşımı, şablon kütüphanesi).

  \item \textbf{B2B Satış:} Hukuk büroları, muhasebe firmaları ve KOBİ dernekleriyle ortaklık;
  çalışan sayısına göre kurumsal plan.

  \item \textbf{API Ekosistemi:} E-imza platformları, ERP/CRM sistemleri ve muhasebe yazılımlarıyla
  entegrasyon; iş ortağı kanalı geliri.

  \item \textbf{Dikey Genişleme:} Faz 2'de gayrimenkul, İK ve finans sektörlerine özel
  playbook şablonları ve dikey SaaS.
\end{enumerate}

% ═════════════════════════════════════════════════════════════════════════════
\section{Rekabet Analizi}
% ═════════════════════════════════════════════════════════════════════════════

\subsection{Rekabet Haritası}

\begin{longtable}{p{3.5cm} >{\centering}p{2cm} >{\centering}p{2cm} >{\centering}p{2.5cm} >{\centering\arraybackslash}p{2.5cm}}
\toprule
\textbf{} & \textbf{Türkçe\newline Destek} & \textbf{Playbook} & \textbf{Çok-Ajanlı\newline AI} & \textbf{KOBİ\newline Odaklı} \\
\midrule
\textbf{\color{primary}Lagent} & \color{success}\textbf{Tam} & \color{success}\textbf{Evet} & \color{success}\textbf{Evet} & \color{success}\textbf{Evet} \\[4pt]
Kira Bul / Hukuk SaaS & \color{success}Evet & \color{danger}Hayır & \color{danger}Hayır & \color{warning}Kısmi \\[4pt]
Harvey AI & \color{danger}İngilizce & \color{success}Evet & \color{success}Evet & \color{danger}Hayır \\[4pt]
ContractPodAi & \color{warning}Kısmi & \color{success}Evet & \color{warning}Kısmi & \color{danger}Hayır \\[4pt]
Spellbook & \color{danger}İngilizce & \color{warning}Kısmi & \color{warning}Kısmi & \color{danger}Hayır \\[4pt]
Genel LLM (ChatGPT) & \color{success}Evet & \color{danger}Hayır & \color{danger}Hayır & \color{danger}Hayır \\
\bottomrule
\end{longtable}

\subsection{Sürdürülebilir Rekabet Avantajı}

\begin{itemize}[itemsep=5pt]
  \item \textbf{Türkçe Yerel Uzmanlık:} Hukuki Türkçe için fine-tune edilmiş prompt mühendisliği
  ve sınıflandırma sistemi. Uluslararası rakipler için dil bariyeri ciddi giriş engeli.

  \item \textbf{Playbook Veri Uçurumu:} Kullanıcılar Playbook kuralları ekledikçe sistem daha
  akıllı hale gelir. Kurumsal müşterilerin geçiş maliyeti (switching cost) artar; ağ etkisi
  çok-kiracılı şablon kütüphanesiyle güçlenir.

  \item \textbf{Çok-Ajanlı Orkestrasyon:} Madde ayrıştırma + risk değerlendirmesi + müzakere
  üretimi zincirlenmesi ve deterministik kural motoru entegrasyonu basit LLM wrapper'larından
  çok daha güvenilir ve özelleştirilebilir çıktı üretir.

  \item \textbf{Denetim İzi + Uyum:} KVKK ve kurumsal uyum gerekliliklerini karşılayan tam
  denetim kaydı; büyük işletmelerde satın alma engelini ortadan kaldırır.

  \item \textbf{Hız:} MVP'den gerçek ürüne kadar 3 aylık geliştirme; teknik borç minimize
  edilmiş modern teknoloji yığını.
\end{itemize}

% ═════════════════════════════════════════════════════════════════════════════
\section{Yol Haritası}
% ═════════════════════════════════════════════════════════════════════════════

\begin{longtable}{>{\bfseries}p{2.5cm} >{\color{primary}\bfseries}p{3cm} p{8cm}}
\toprule
\textbf{Dönem} & \textbf{Aşama} & \textbf{Hedefler} \\
\midrule
\endhead
Q2 2026 & \color{success}MVP (Tamamlandı) &
Madde ayrıştırma, risk değerlendirmesi, müzakere ajanı, onay akışı, raporlama.
Kullanıcı kaydı ve kimlik doğrulama sistemi. \\[6pt]

Q3 2026 & Kapalı Beta &
500 beta kullanıcısı; hukuk bürosu ve KOBİ ortaklıkları.
Kullanıcı geri bildirimine göre UX iyileştirmesi.
OCR kalitesi ve sınıflandırma doğruluğu artırımı.
Kira, tedarik ve istihdam için sektör playbook şablonları. \\[6pt]

Q4 2026 & Genel Çıkış &
Ücretli abonelik lansmanı (Starter + Pro).
API açılımı ve ilk entegrasyon ortakları.
Hedef: 2.000 ücretli kullanıcı, 2M TL ARR. \\[6pt]

Q1–Q2 2027 & Büyüme &
Enterprise satış ekibi.
E-imza platformları entegrasyonu (e-imza.com.tr, DocuSign).
Çok-dilli destek (İngilizce, Almanca).
Hedef: 15.000 ücretli kullanıcı, 15M TL ARR. \\[6pt]

2027–2028 & Ölçekleme &
Orta Doğu ve DACH pazarlarına açılım.
Dikey SaaS (gayrimenkul, İK, finans).
White-label kurumsal çözümler.
Hedef: 100.000 kullanıcı, 100M TL ARR. \\
\bottomrule
\end{longtable}

% ═════════════════════════════════════════════════════════════════════════════
\section{Ekip}
% ═════════════════════════════════════════════════════════════════════════════

\begin{highlightbox}
Kurucu ekip teknik yetkinlik ve ürün vizyonunu birleştiriyor. Üç kurucu da projeyi sıfırdan
inşa etmiş; çok-ajanlı AI mimarisi, hukuki iş akışları ve modern web geliştirme konusunda
derinlemesine deneyim kazanmıştır.
\end{highlightbox}

\vspace{0.5em}

\begin{longtable}{>{\bfseries}p{4.5cm} p{3.5cm} p{6cm}}
\toprule
\textbf{Ad} & \textbf{Rol} & \textbf{Uzmanlık} \\
\midrule
Mustafa Emir Ceran & Kurucu CEO / Teknik Lider &
Çok-ajanlı AI orkestrasyon, LLM entegrasyonu, LangGraph, ürün vizyonu ve proje yönetimi \\[6pt]
Mert Ayrancı & Kurucu CTO / Backend &
FastAPI, PostgreSQL, RAG mimarisi, risk ajanı geliştirme, API tasarımı \\[6pt]
Osman Gazi Atalay & Kurucu CPO / Frontend &
Next.js, React, kullanıcı deneyimi, redline görüntüleyici, QA ve test otomasyonu \\
\bottomrule
\end{longtable}

\vspace{0.5em}
\textbf{Danışman İhtiyaçları:} Hukuk sektörü domain uzmanı (avukat / hukuk teknolojileri), B2B
satış deneyimli büyüme danışmanı ve fintech/legaltech yatırım geçmişi olan stratejik yatırımcı.

% ═════════════════════════════════════════════════════════════════════════════
\section{Finansal Projeksiyonlar}
% ═════════════════════════════════════════════════════════════════════════════

\subsection{Gelir Projeksiyonu (3 Yıl)}

\begin{center}
\begin{tikzpicture}
\begin{axis}[
  ybar,
  bar width=1.5cm,
  width=12cm, height=6cm,
  symbolic x coords={2026 (Q4), 2027, 2028},
  xtick=data,
  ylabel={TL (Milyon)},
  ymin=0, ymax=120,
  ymajorgrids=true,
  grid style={dashed, gray!30},
  every axis plot/.append style={fill=accent, fill opacity=0.85, draw=primary},
  legend style={at={(0.5,-0.2)}, anchor=north, legend columns=-1},
  tick label style={color=dark},
  label style={color=dark},
  axis line style={color=muted},
]
\addplot coordinates {(2026 (Q4), 2) (2027, 15) (2028, 80)};
\end{axis}
\end{tikzpicture}
\end{center}

\begin{longtable}{p{4cm} r r r}
\toprule
\textbf{Metrik} & \textbf{2026 (Q4)} & \textbf{2027} & \textbf{2028} \\
\midrule
Ücretli Kullanıcı & 2.000 & 15.000 & 75.000 \\
ARR (TL Milyon) & 2 & 15 & 80 \\
Brüt Marj & \%65 & \%70 & \%75 \\
Müşteri Edinme Maliyeti & 1.500 TL & 1.200 TL & 900 TL \\
Aylık Büyüme Oranı & \%25 & \%18 & \%12 \\
\bottomrule
\end{longtable}

\subsection{Maliyet Yapısı}

\begin{itemize}[itemsep=3pt]
  \item \textbf{LLM API Maliyeti:} Sözleşme başına $\sim$0.8–1.5 TL (GPT-4o-mini kullanımında);
  hacim artışıyla azalır; önbellekleme ile \%30–40 tasarruf sağlanabilir.
  \item \textbf{Altyapı:} Cloud hosting + veritabanı $\sim$15.000 TL/ay (MVP); ölçeklendikçe
  aylık 80.000–150.000 TL.
  \item \textbf{İnsan Kaynağı:} 3 kurucu + 2 mühendis (Faz 2); büyüme turunda satış + müşteri
  başarısı kadrosu.
\end{itemize}

% ═════════════════════════════════════════════════════════════════════════════
\section{Yatırım Talebi}
% ═════════════════════════════════════════════════════════════════════════════

\begin{statbox}
\begin{center}
{\color{primary}\Large\textbf{Hedef Yatırım Tutarı: 5.000.000 TL}}\\[0.5em]
{\color{muted}Seri A Öncesi Tur — Mayıs 2026}\\[1em]
\begin{tabular}{lrl}
\textbf{Ürün Geliştirme} & \%35 & 1.750.000 TL \\
\textbf{Pazarlama ve Büyüme} & \%30 & 1.500.000 TL \\
\textbf{İnsan Kaynağı} & \%25 & 1.250.000 TL \\
\textbf{Altyapı ve Operasyon} & \%10 & 500.000 TL \\
\end{tabular}
\end{center}
\end{statbox}

\subsection{Kullanım Detayı}

\subsubsection{Ürün Geliştirme (1.750.000 TL)}
\begin{itemize}[itemsep=3pt]
  \item Sınıflandırma doğruluğunu \%70 → \%90'a çıkaracak model fine-tuning ve prompt
  optimizasyonu
  \item Mobil uygulama (iOS / Android) geliştirmesi
  \item API ürünü ve geliştirici ekosistemi altyapısı
  \item E-imza ve ERP/CRM entegrasyonları
\end{itemize}

\subsubsection{Pazarlama ve Büyüme (1.500.000 TL)}
\begin{itemize}[itemsep=3pt]
  \item Dijital pazarlama (SEO, içerik, sosyal medya)
  \item Hukuk bürosu ve KOBİ birlikleriyle ortaklık programı
  \item Beta kullanıcı kampanyası ve referans programı
  \item Etkinlik katılımı ve marka bilinirliği
\end{itemize}

\subsubsection{İnsan Kaynağı (1.250.000 TL)}
\begin{itemize}[itemsep=3pt]
  \item 2 senior backend/AI mühendisi
  \item 1 ürün tasarımcısı (UX/UI)
  \item 1 B2B satış yöneticisi
\end{itemize}

\subsection{Beklenen Çıktılar (18 Ay)}

\begin{itemize}[itemsep=5pt]
  \item \textbf{15.000 ücretli kullanıcıya} ulaşmak
  \item \textbf{15M TL ARR} gerçekleştirmek
  \item \textbf{Seri A hazırlığı:} 30–50M TL değerleme hedefli Seri A turu için metriklere ulaşmak
  \item \textbf{3 kurumsal müşteri} (hukuk bürosu veya büyük KOBİ) sözleşmesi imzalamak
\end{itemize}

% ═════════════════════════════════════════════════════════════════════════════
\section*{Kapanış: Neden Lagent, Neden Şimdi?}
% ═════════════════════════════════════════════════════════════════════════════

\begin{highlightbox}
\begin{enumerate}[leftmargin=*, itemsep=6pt]
  \item \textbf{Zamanlama:} GPT-4 kalitesindeki LLM'lerin maliyeti 2023–2026 arasında \%90
  düştü. Türkçe NLP olgunlaştı. Pazar hazır.

  \item \textbf{Problem-Çözüm Uyumu:} Sözleşme analizi sınırlı, tekrarlanabilir ve ölçülebilir
  bir görevdir. AI bu görev için ideal; halüsinasyon riski Playbook + kural motoru ile kontrol
  altına alınmıştır.

  \item \textbf{Savunulabilir Avantaj:} Türkçe hukuki domain uzmanlığı + Playbook veri uçurumu +
  çok-ajanlı mimari rakiplerin kısa vadede eşleşemeyeceği bir kombinasyon oluşturuyor.

  \item \textbf{Devasa Pazar:} 8M+ KOBİ ve 85M nüfus; dijital sözleşme hacmi yılda \%40 büyüyor.

  \item \textbf{Çalışan Ürün:} Sıfırdan inşa edilmiş, gerçek kullanıcılarla test edilmiş MVP.
  Sadece fikir değil — teslim etme kapasitesi kanıtlanmış bir ekip.
\end{enumerate}
\end{highlightbox}

\vspace{1.5em}

\begin{center}
{\color{primary}\Large\textbf{Birlikte Türkiye'nin Hukuki Altyapısını Dönüştürelim.}}\\[1em]
{\color{muted}\normalsize persevia.ai@gmail.com}
\end{center}

\vspace{2em}
\begin{center}
\rule{8cm}{0.4pt}\\[0.5em]
{\small\color{muted}Bu belge yatırımcı bilgilendirme amaçlıdır. Finansal projeksiyon ve piyasa
tahminleri mevcut verilere dayalı öngörüler içermekte olup garanti niteliği taşımamaktadır.}
\end{center}

\end{document}
